\section{Discussion}\label{sec:discussion}
The compact A3 model demonstrates parameter reduction and an observed memory gain against E5 on these fixed new services. The broader evidence does not support a uniform accuracy improvement. CPU A3 is worse in mean MASE than both primary references, the A0 variant has lower memory MASE than A3, and the strongest mean MASE in the primary cohort belongs to the small PatchTST recipe. A3's favorable memory contrast is concentrated in a minority of services, although its seed direction is consistent under the prespecified fixed-cohort comparison. The outcome-informed service exclusions show why directional seed consistency alone does not imply broad spatial portability.

The algebraic result explains a practical boundary: a fixed MA decomposition followed by separate affine heads is itself affine. Reducing its parameterization can lower measured model-forward cost, but a different Adam trajectory and explicit availability coefficients can improve or worsen errors. E5's availability-dependent EMA is outside that exact collapse. The independently tuned A1/A2/A3 means do not demonstrate a general additive benefit from either simple missingness descriptor. A future mechanism study could hold hyperparameters fixed, evaluate descriptor interactions against matched regularization and use a longer independently sampled period; those are proposed experiments, not results here.

The external case changes both data domain and measurement definition. Its persistence advantage is consistent with a short, slowly changing normalized trace, but these observations do not establish the underlying cause. No external graph was available, and six calibration hours yield potentially small scale denominators. We therefore report external MASE and normalized absolute errors together, without combining their magnitudes with ChronoGraph physical errors. Additional days, independently chosen services and datasets with documented preprocessing would test whether any ranking persists.

The efficiency measurements add a useful but bounded perspective. A3 is much smaller than E5, and its resident-tensor p50 latency is lower on both measured devices. PatchTST offers lower primary mean MASE at a higher measured forward cost; simple moving average is a strong CPU accuracy reference with no trainable weights. In a small batch, CPU sometimes outperforms CUDA because launch and synchronization costs matter. None of these facts establishes lower energy, quicker telemetry collection or better operational autoscaling. A deployment trial would need a controller, workload replay, end-to-end timing, explicit SLA metrics and a reliable project-level power meter.

The preserved historical study should be read with its separate condition definitions. Its E5 memory ranking was a hypothesis source, not independent confirmation. The new spatial comparison demonstrates a material change in model ordering, and Alibaba changes it again. A severe historical GinAR error is retained as an implementation/protocol result and does not warrant a claim about the model's published behavior. The principal reproducible contribution is the transparent accounting of when compactness, error, missingness and domain transfer disagree.


The historical-to-expanded transition changes more than the number of seeds. The destination services have different per-service calibration, a different role in fitting, and a source training recipe that now includes ten conditions. E5's change in rank consequently cannot be attributed solely to having more services or initializations. Historical memory improvement, C-family failure to preserve it, and the new A3 contrast jointly answer the compactness question: reducing coefficients can retain an acceptable measured error while losing a specific earlier advantage, and a later favorable reference comparison does not restore universal superiority.

NLinear and PatchTST suggest that an explicit recent level or a learned patched representation can be competitive under source-trained transfer. The experiment does not isolate those mechanisms. Their selected optimization settings and parameterizations differ, and the PatchTST recipe contains normalization and attention together. Moreover, NLinear has slightly lower displayed CPU MAE than PatchTST despite the opposite MASE order. MASE weights errors relative to each service's training variability, whereas absolute macro errors weight the same physical unit. An algorithm ranking therefore depends on the scientific loss function as well as the cohort.

Longer direct horizons increase displayed errors for the selected learned models, but these endpoints use independently fitted weights and common target support. They are not evidence of recursive-error accumulation. Longer histories have mixed resource-specific effects: improved CPU A3 and worsened memory A3 are both retained. The counterexample matters because a longer window increases coefficients, input processing and possible exposure to older regimes. It cannot be assumed to improve generalization merely by containing more values. The star design describes these particular choices rather than estimating a complete interaction between every history and horizon.

Finally, the historical correction of interval units is scientifically substantive. An absolute MASE interval cannot be read as a percentage interval simply because it appears next to a percentage effect. In the expanded analysis, mean paired cell reductions and ratios of means remain explicitly separate. Likewise, the sign test measures seeded directional consistency and the conditional block interval describes a resampled mean. Agreement is informative, but disagreement is possible without an arithmetic error. These distinctions prevent narrow conditional uncertainty or a favorable rank from becoming an unsupported population-wide superiority claim.

The source training-only dynamics provide one check on the resource-specific interpretation. Median lag-one autocorrelation across the 28 series is 0.447 for CPU and 0.867 for memory. Median absolute change relative to series standard deviation is 0.413 and 0.275, respectively. This is compatible with more persistent memory levels. It cannot identify why the EMA changes forecast error, because the residual retains deviations from the estimated level and the two heads are fitted jointly. The observations motivate a controlled error analysis around abrupt changes rather than a claim that smoothing mechanically delays the final forecast.
